\section{Introduction}
\label{sec:intro}

Software maintenance accounts for the majority of the total cost of the software
development life cycle, and locating and fixing defects is among its most effort-intensive
activities~\cite{}. Automated Program Repair (APR) aims to reduce this
burden by generating candidate fixes for defects automatically~\cite{legoues2019apr}, and
it has entered an agentic era~\cite{}. Rather than producing a
patch in a single step, modern repair systems built on Large Language Models (LLMs)
operate as autonomous agents that explore a repository, localize the fault, edit the
code, and run the project's tests before proposing a fix~\cite{yang2024sweagent,
zhang2024autocoderover}. Such tools are increasingly adopted in software development
practice~\cite{so2024, so2025}. At the same time, benchmarks built from real-world software
issues have advanced their capability: on SWE-bench, a suite of bugs drawn from
open-source \textsc{GitHub} projects~\cite{jimenez2024swebench}, and its human-validated
subset \sweBenchVerified{}~\cite{swebenchverified2024}, the fraction of issues that
agents resolve has risen within a few years.

However, this adoption has not been matched by a corresponding growth in trust.
According to the Stack Overflow Developer Survey, the share of developers using or
planning to use AI tools rose from \aiAdoptionPrev{} in 2024 to \aiAdoptionCur{} in
2025~\cite{so2024, so2025}. Over the same period, however, developers' trust in the
\accuracy{} of AI output moved in the opposite direction, with the proportion who
trust it falling from \aiTrustPrev{} to \aiTrustCur{} while active distrust rose from
\aiDistrustPrev{} to \aiDistrustCur{}~\cite{so2024, so2025}. This widening gap between
reliance and confidence is a symptom of the miscalibrated trust that prior work
identifies as a central risk of adopting AI for software engineering~\cite{khati2025mapping}.
Moreover, trust is not a single quantity. In a recent study of trust in LLM-based
software engineering, practitioners ranked \accuracy{}, \interpretability{}, and \robustness{} as
the antecedents that matter most for program repair~\cite{khati2025mapping}. The same study
distinguishes
\emph{trust}, a subjective user attitude, from \emph{trustworthiness}, an intrinsic and
measurable property of the system. It is the latter that remains largely unmeasured
for agentic program repair.

These perceptions point to an underlying technical problem. In the same survey,
\aiAlmostRight{} of developers report encountering AI solutions that are ``almost right,
but not quite''~\cite{so2025}: output that looks correct yet fails on closer inspection. In program
repair, this failure mode is built into how success is measured. The field evaluates a
repair against a single criterion: a patch is considered correct if it passes the
project's test suite. But passing the test suite is only a proxy for a correct fix, and a
partial one: because the suite is necessarily incomplete, a patch can satisfy every
available test while modifying the wrong code and masking, rather than fixing, the
underlying fault. This is the well-studied problem of patch
overfitting~\cite{smith2015overfitting, qi2015analysis}. It is not a hypothetical concern
for agentic repair: even on the human-validated \sweBenchVerified{}, recent studies
find that between \divergeLow{} and \divergeHigh{} of test-passing patches diverge from the
developer's ground-truth fix, with roughly \certIncorrect{} outright
incorrect~\cite{wang2026solved, yu2025utboost}. A developer therefore cannot safely trust a
test-passing patch without understanding \emph{why} it works. We argue that
trustworthiness should instead be measured directly, by recasting the three antecedents
that practitioners rank as mattering most for program repair (\accuracy{}, \interpretability{},
and \robustness{}) as concrete, checkable properties of a repair~\cite{khati2025mapping}.
\accuracy{} asks whether the patch actually fixes the bug. \interpretability{}
asks whether the agent repairs the bug for the right reasons, and decomposes into
\emph{localization validity} (the agent identifies the true fault location before editing)
and \emph{rationale alignment} (the agent's stated reasoning correctly identifies the root
cause of the bug). \robustness{} asks whether \accuracy{} and \interpretability{} persist
when the bug is semantically transformed. Together, these axes expose a failure that a
pass/fail bar cannot: an agent may resolve a bug with a patch whose localization misses
where the fault lies, fixing the bug without pointing at it, a pattern we call
\emph{fix-without-localizing}. This pattern is not hypothetical: in our study, \mislocRate{}
of the patches that pass the test suite are localized to the wrong line, and \wrongFileRate{}
do not edit the faulty file at all.

Evaluating these axes, however, is not straightforward. \accuracy{} is readily measured by
executing the project's test suite, but \interpretability{} and \robustness{} are not.
Assessing \interpretability{} requires, for every bug, a trusted ground truth for both the
fault location and the root-cause mechanism, together with a way to compare an agent's
free-form natural-language rationale against that mechanism at scale. Human judgment
provides the reference here, but it does not scale to thousands of agent runs, whereas an
automated judge must itself be shown to agree with humans before its verdicts can be
trusted. Assessing robustness, in turn, requires transformations that alter a program's
surface form while preserving its semantics, so that any change in an agent's behavior can
be attributed to the transformation rather than to a change in the task. Meeting these
three requirements is the methodological challenge we address.

To this end, we introduce \trustRepair{}, a benchmark and evaluation pipeline that
measures the trustworthiness of agentic program repair along all three axes.
\trustRepair{} builds on \numBugs{} real-world bugs from \sweBenchVerified{}, for which two annotators independently established the ground-truth fault
locations and root-cause rationales. To measure rationale alignment at scale, we employ a
panel, or \emph{jury}, of three LLMs~\cite{verga2024poll} that compares an agent's reasoning
against the human rationale under order-swapped, identity-blinded prompts, following best
practices for LLM-based evaluation~\cite{zheng2023judging}; we validate this jury against
the human labels and find that it agrees with them \juryAgreement{} of the time. To measure
robustness, we apply \numTransforms{} semantics-preserving transformation rules to each bug and
re-evaluate \accuracy{} and \interpretability{} on the transformed variant. Using this pipeline,
we conduct an empirical study of four recent repair agents, each run three times
per bug, and investigate two research questions:

\begin{itemize}
  \item \textbf{RQ1 (Trust axes and their alignment).} How well do these
  agents perform on \accuracy{} and \interpretability{}, and do these axes agree at the level of
  individual bugs, or do agents fix bugs without localizing or understanding them?
  \item \textbf{RQ2 (Robustness of trust).} Do \accuracy{} and \interpretability{} persist when a
  bug is subjected to semantics-preserving transformations, and which axes degrade most?
\end{itemize}

Our study yields several findings that a single-axis, \accuracy{}-only evaluation would
obscure. First, single-axis competence substantially overstates trustworthiness: even the
strongest agent earns \emph{full trust} (a correct fix that is both correctly localized
and correctly explained) on only \fullTrustBest{} of bugs, well below its \accuracy{}. Second,
fix-without-localizing is common rather than exceptional: \mislocRate{} of the patches that
pass the test suite are localized to the wrong line and \wrongFileRate{} miss the faulty file
entirely. This gap is one of localization rather than comprehension: \understoodShare{} of
these mislocalized-but-resolved patches still describe the bug's mechanism correctly, so
localization, not rationale, is the weaker link in the interpretability chain. Third, the
axes are only weakly coupled at the instance level, so an agent that ranks well on one axis
routinely fails another on the same bug. \textcolor{red}{[TODO: will be added later]} In summary, this paper makes the following contributions:

\begin{itemize}
  \item We introduce a multi-axis \emph{trustworthiness} benchmark and a reusable,
  human-validated evaluation pipeline for agentic program repair that jointly measures
  \accuracy{}, \interpretability{}, and \robustness{}.
  \item We conduct an empirical study of four recent repair agents that reveals a
  systematic dissociation between resolving a bug and localizing it: \mislocRate{} of
  successful repairs are localized to the wrong line, most of them while still explaining the
  bug correctly.
  \item We release our benchmark, ground-truth annotations, jury, and leaderboard to
  support replication and future research.\footnote{\textcolor{red}{[TODO: anonymized
  artifact link for the double-anonymous submission.]}}
\end{itemize}

The remainder of this paper is organized as follows. Section~\ref{sec:related} surveys
related work; Section~\ref{sec:problem} defines the trust axes and research questions;
Section~\ref{sec:benchmark} describes the construction of \trustRepair{};
Section~\ref{sec:method} details our measurement methodology; Section~\ref{sec:results}
reports our results; Sections~\ref{sec:discussion} and~\ref{sec:threats} discuss
implications and threats to validity; and Section~\ref{sec:conclusion} concludes.
